\section{Limitations}
\label{sec:limitations}

First, all behavior is simulated. Hardware contact uncertainty and natural variation in human instructions remain untested. The causal object-layout conclusions are restricted to the positive-control-passing DROID/RoboLab rows; RoboTwin screening results cannot be read as cross-arena replication or failure.

Second, the task family contains two named objects, one static instruction, and one lateral relation. C002 tests one controlled reference inversion for one $\pi_{0.5}$ checkpoint in one symmetric scene. It is not a broad paraphrase benchmark and does not establish unrestricted compositional grounding.

Third, reflection moves object centers, whereas the symmetric condition changes a registered scene package that also includes companion-object inventory. Both leave embodiment, camera viewpoint, wrist mounting, action parameterization, and decoder unchanged. They identify dependence of the measured outcome on scene configuration; they do not localize errors inside the model or robot system. Checkpoint provenance is also not fully auditable for every evaluated model, so we avoid architecture-level attribution.

Finally, the fixed-observation diagnostic is action-space only and does not bind prompt-conditioned actions to verified forward kinematics. These limits narrow the causal interpretation, but they do not alter the central evaluation result: the same binary directional score can change under reference inversion or scene interventions.
